\subsection{Support arcs and angular curvature without boundary regularity}
\label{subsec:nonsmooth-support-curvature}
The incoming arcs can be treated through sections, without a
parametrization by the boundary normal.  This also distinguishes
strict convexity from smoothness in the subsequent measure
calculation.

\begin{lemma}[A section form of the arc identity]
\label{lem:nonsmooth-arc-support}
Let $C$ be a strictly convex planar body with nonempty interior.
Let $p_-,p_+$ be its unique contacts with the support lines of
normals $-e^\perp,e^\perp$, and put $L=[p_-,p_+]$.
For its closed incoming arcs $\Gamma_+,\Gamma_-$ one has
\begin{align}
 h_{\Gamma_+}(u)&=
 \begin{cases}h_C(u),&u\cdot e\le0,\\
                h_L(u),&u\cdot e\ge0,
 \end{cases}
 \label{eq:nonsmooth-plus-hemispheres}\\
 h_{\Gamma_-}(u)&=
 \begin{cases}h_L(u),&u\cdot e\le0,\\
                h_C(u),&u\cdot e\ge0.
 \end{cases}
 \label{eq:nonsmooth-minus-hemispheres}
\end{align}
The two alternatives agree when $u\cdot e=0$.  In particular,
\begin{equation}\label{eq:nonsmooth-arc-identity}
                  h_{\Gamma_+}+h_{\Gamma_-}=h_C+h_L.
\end{equation}
\end{lemma}
\begin{proof}
Write $u=\alpha e+\beta e^\perp$ and use the interval sections
$[\ell(w),r(w)]$ for $w\in[w_-,w_+]$.  Convexity of $C$ makes
$\ell$ convex and $r$ concave.  They are continuous on the
closed projection interval: interior continuity follows from
convexity, and endpoint continuity follows from compactness and
the uniqueness of the extreme sections.  The endpoints of both
graphs are $p_-,p_+$.

If $\alpha\le0$, maximizing $\alpha s+\beta w$ over each section
chooses $s=\ell(w)$.  Consequently
$h_{\Gamma_+}(u)=h_C(u)$.  If $\alpha\ge0$, the function
$w\mapsto\alpha\ell(w)+\beta w$ is convex.  Its value at an
interior point is at most the larger endpoint value, by the
defining convexity inequality.  Its maximum is therefore
$\max\{u\cdot p_-,u\cdot p_+\}=h_L(u)$.
For the other arc, $\alpha\ge0$ selects the right endpoint of
each section and gives $h_C$; when $\alpha\le0$ the function
$\alpha r(w)+\beta w$ is convex and gives $h_L$.
At $\alpha=0$ both maxima depend only on the projection interval.
This proves all the formulas, including directions supporting
a corner and the two transverse directions.
\end{proof}

We identify the angular circle with $\R/(2\pi\Z)$, put
$n_\phi=(\cos\phi,\sin\phi)$ and
$\tau_\phi=n_\phi^\perp$, and use the planar support convention
\begin{equation}\label{eq:nonsmooth-area-convention}
 \int\psi\dd\mathsf S_K
     =\int_0^{2\pi}h_K(n_\phi)
                      (\psi''(\phi)+\psi(\phi))\dd\phi,
 \qquad \psi\in C^\infty(\R/(2\pi\Z)).
\end{equation}
Thus angular differentiation is distributional.  The following
proof supplies both positivity and the precise atom convention
for this surface-area measure.

\begin{lemma}[Surface-area atoms and exposed segments]
\label{lem:nonsmooth-surface-atoms}
For a compact convex planar body $K$ with nonempty interior,
\eqref{eq:nonsmooth-area-convention} defines a finite nonnegative
measure.  Let
\[
 F(K,n_\phi)=\{x\in K:x\cdot n_\phi=h_K(n_\phi)\}
\]
be its exposed face.  Writing $h(\phi)=h_K(n_\phi)$, its
one-sided derivatives satisfy
\begin{equation}\label{eq:nonsmooth-face-derivatives}
 h'_-(\phi)=\min_{x\in F(K,n_\phi)}x\cdot\tau_\phi,
 \qquad
 h'_+(\phi)=\max_{x\in F(K,n_\phi)}x\cdot\tau_\phi.
\end{equation}
Consequently
\begin{equation}\label{eq:nonsmooth-atom-face-length}
 \mathsf S_K(\{n_\phi\})
       =h'_+(\phi)-h'_-(\phi)
       =\operatorname{length}F(K,n_\phi).
\end{equation}
In particular, $\mathsf S_K$ is nonatomic if and only if $K$ is
strictly convex.  Boundary corners and a singular continuous
part of $\mathsf S_K$ are permitted in this equivalence.
\end{lemma}
\begin{proof}
Choose $R$ with $K\subset B(0,R)$.  The support function is
Lipschitz in $\phi$.  Each function
$x\cdot n_\phi+R\phi^2/2$, $x\in K$, is convex on the real
line, since its second derivative is $R-x\cdot n_\phi\ge0$.
Their supremum $h(\phi)+R\phi^2/2$ is convex.  Hence $h$ has
one-sided derivatives of locally bounded variation and its
distributional second derivative is a locally finite signed
measure.  Periodicity gives a finite signed measure on the circle.

If $p$ is a support point at $\phi$, then, for $|b|<\pi/2$,
\[
 h(\phi+b)+h(\phi-b)
 \ge p\cdot(n_{\phi+b}+n_{\phi-b})
 =2\cos b\,h(\phi).
\]
Multiply by a nonnegative smooth periodic $\psi$, integrate,
translate the two shifted integrals, and divide by $b^2$.
Letting $b\to0$ yields
$\int h(\psi''+\psi)\ge0$.  Thus the signed measure
$h''+h\dd\phi$ is nonnegative and is the measure in
\eqref{eq:nonsmooth-area-convention}.

For $b>0$ and $p\in F(K,n_\phi)$, the support inequality gives
\[
 \frac{h(\phi+b)-h(\phi)}{b}
       \ge p\cdot\frac{n_{\phi+b}-n_\phi}{b}.
\]
Conversely, choose $p_b\in F(K,n_{\phi+b})$.  The same
quotient is at most
$p_b\cdot(n_{\phi+b}-n_\phi)/b$.
Every subsequential limit of $p_b$ belongs to $F(K,n_\phi)$,
by compactness and continuity of $h$.  Taking lower and upper
limits proves the right-derivative formula in
\eqref{eq:nonsmooth-face-derivatives}; negative $b$ gives the
left-derivative formula.  The atom of the Stieltjes measure
$h''$ at $\phi$ is the jump $h'_+(\phi)-h'_-(\phi)$, whereas
$h\dd\phi$ has no atoms.  The face is a segment on the line
$x\cdot n_\phi=h(\phi)$, so its length is exactly the difference
between the two extremal $\tau_\phi$ coordinates.  This proves
\eqref{eq:nonsmooth-atom-face-length}.

The geometric interpretation of the measure follows in the same
coordinates.  The right-continuous support contact
$x_+(\phi)=h(\phi)n_\phi+h'_+(\phi)\tau_\phi$ selects the
positive tangent endpoint of each exposed face.  As the normal
turns once, these contacts traverse the convex boundary in order;
each jump is completed by the corresponding exposed segment.
The distributional product rule gives
\[
                      D x_+=\tau_\phi\dd\mathsf S_K.
\]
Since $|\tau_\phi|=1$ and $\mathsf S_K\ge0$, the variation
of this boundary parametrization is $\mathsf S_K$.  It assigns
boundary arclength to its outer normal, with a whole exposed
segment assigned to its fixed normal.  This is the usual planar
surface-area measure in the convention
\eqref{eq:nonsmooth-area-convention}.

A strictly convex body has only singleton exposed faces and
hence no atoms.  Conversely, any nontrivial boundary segment
lies in an exposed face: a supporting line at an interior point
of the segment contains both endpoints.  Such a segment would
give a positive atom.  At a corner of a strictly convex body,
the normal may vary over an interval while the support point
remains a fixed point $p$.  On the interior of that interval
$h(\phi)=p\cdot n_\phi$, so $h''+h=0$ there; its endpoints
also have zero atomic mass by
\eqref{eq:nonsmooth-atom-face-length}.  None of these arguments
excludes a singular continuous part of the measure.
\end{proof}

These are the planar support-measure conventions of
\cite[Chapter~4]{Schneider}.  Signed support differences and the
separate atomic, absolutely continuous and singular continuous
parts of their length measures are treated in
\cite[Section~4]{MartinezMaure}.  Formula
\eqref{eq:nonsmooth-area-convention} retains all three parts;
an almost-everywhere second derivative alone would not do so.

\begin{lemma}[Jordan recovery of the contact chord]
\label{lem:nonsmooth-jordan-chord}
For the body and chord in Lemma~\ref{lem:nonsmooth-arc-support},
put $H=h_C-h_L$, $m=(p_-+p_+)/2$, and
$\ell=|p_+-p_-|$.  If $n$ is either unit normal of the chord,
then the Jordan decomposition of $\sigma_H=(\partial_\phi^2+1)H$
is
\begin{equation}\label{eq:nonsmooth-jordan-parts}
 \sigma_H^+=\mathsf S_C,\qquad
 \sigma_H^-=\ell(\delta_n+\delta_{-n}).
\end{equation}
It determines the centered chord
$L^0=[-\ell n^\perp/2,\ell n^\perp/2]$, and
\begin{equation}\label{eq:nonsmooth-jordan-recovery}
                     H+h_{L^0}=h_{C-m}.
\end{equation}
Exactly one of the two chord normals has positive scalar product
with $e$.
\end{lemma}
\begin{proof}
The projection interval has positive length, so $\ell>0$.
Writing $q=(p_+-p_-)/\ell$, the chord support is
\[
 h_L(n_\phi)=m\cdot n_\phi+\frac\ell2|q\cdot n_\phi|.
\]
The translation term is annihilated by $\partial_\phi^2+1$.
The second term solves $h''+h=0$ away from the two normals
$n,-n$, and its first derivative jumps upward by $\ell$ at
each.  Hence
\[
 \sigma_H=\mathsf S_C-\ell(\delta_n+\delta_{-n}).
\]
Lemma~\ref{lem:nonsmooth-surface-atoms} shows that $\mathsf S_C$
assigns zero mass to these two points.  The two nonnegative
measures on the right are therefore mutually singular, so they
are exactly the positive and negative Jordan parts.  The
locations and equal masses of the negative atoms determine
$L^0$, independent of the choice between $n$ and $-n$.
Since $h_L=h_{L^0}+m\cdot u$, adding $h_{L^0}$ to $H$ gives
\eqref{eq:nonsmooth-jordan-recovery}.  Finally,
$(p_+-p_-)\cdot e^\perp=w_+-w_->0$; a chord normal therefore
has nonzero scalar product with $e$, and its sign selects exactly
one of the two normals.  The chord need not be perpendicular to
the command direction.
\end{proof}
